\providecommand{\BM}{\textbf{[B]}}
\providecommand{\KM}{\textbf{[K]}}
\providecommand{\SM}{\textbf{[S]}}
\providecommand{\XM}{\textbf{[X]}}
\providecommand{\QM}{\textbf{[Q]}}

\chapter*{Interfaces of FFGFT to Other Frameworks}
\addcontentsline{toc}{chapter}{Interfaces of FFGFT to Other Frameworks}

{\small\noindent\textbf{Abstract.}\quad\ignorespaces
The corpus contains comparison and bridge documents for eleven frameworks,
spread over more than twenty documents. This document collects them: for each
framework the document numbers, the mathematical point of contact, the status
and the open item. The quantities where several frameworks actually overlap
follow, each with its route, as a template for a comparison table. Nothing new
is computed; all numbers come from the documents named and are recomputed in
the check script.
}\medskip

\section*{List of symbols}
\addcontentsline{toc}{section}{List of symbols}

\medskip
\begin{tabular}{@{}ll@{}}
\textbf{[B]} & algebraically proven \\
\textbf{[K]} & numerically confirmed \\
\textbf{[S]} & structurally plausible, derivation open \\
\textbf{[X]} & checked and negative \\
\textbf{[Q]} & taken from an external source, cited \\
\end{tabular}

\section{Principle}

FFGFT derives from $\xi$. Measured values and Standard Model quantities serve
for conversion and testability, not for derivation; the cosmic sector stays
bracketed (P39). Quantities of external frameworks appear in all documents
named below as cited inputs \QM{} and do not stand on equal footing with our
own posits. A point of contact counts as a bridge only when the same
structural object follows from each framework's own geometry and an explicit
map between the derivations exists; the criterion is from Doc.~283.

\section{Overview}

\begin{center}
\small
\begin{tabular}{@{}llp{5.2cm}l@{}}
\toprule
\textbf{Framework} & \textbf{Documents} & \textbf{Point of contact} & \textbf{Status} \\
\midrule
HLV (Krüger) & 271, 272, 276, 282, 283, 285, 294, 297 &
  icosahedral cut-and-project quasicrystal versus $T^4/\mathbb{Z}_3$;
  dimension bridge $6=\mathbf{3}\oplus\mathbf{3'}$ in $A_5$; CP divisibility &
  bridge \KM, fork at the symmetry order \\
OPH (Mueller) & 364, 374, 375, 376 &
  spectrum of the icosahedral carrier; energy--length bridge
  $E_{\rm bit}=\hbar c/L$; $v/E_P$; coefficient $8\pi$ &
  bridge equations \BM, scales open on the other side \QM \\
Hyperbit, $G(6,\mathbb{Z}_3\mathbb{C})$ (Matzke) & 324, 326, 336 &
  $\mathbb{Z}_3$ structure and $\mathrm{GF}(9)$; bit energy and the Landauer
  threshold &
  algebraic bridge \BM, bit energy system-dependent \KM \\
XYLATIC (Phillips) & 252 &
  number-theoretic structures of the $\sigma$ function &
  structural comparison \\
RA/PMT (Guevara) & 269 &
  three-layer scheme, perception level &
  structural comparison \\
Vopson & 251 &
  computational-universe hypothesis, second law of infodynamics &
  different primitives \\
Matsas et al. & 105 &
  the number of fundamental constants &
  classification \\
$\Lambda$CDM & 309 &
  scale anchor: both frameworks need exactly one number for the SI coupling &
  shared situation \KM \\
Kagome lattice & 361 &
  four structural parallels; mass as geometric freezing &
  analogy \SM \\
SM Lagrangian & 049 &
  number of fields and free parameters against one equation &
  comparison of form \\
Literature in general & 345 &
  classification against published approaches &
  classification \\
\bottomrule
\end{tabular}
\end{center}

UIFT (Teker) is touched in Docs.~013 and 257 through the bit mass, but has no
comparison document of its own yet.

\section{Quantities where frameworks overlap}

This table is the template for a comparison table. Every value comes with its
route; without it the value is not comparable (Doc.~375).

\begin{center}
\small
\begin{tabular}{@{}llrl@{}}
\toprule
\textbf{Quantity} & \textbf{FFGFT value} & \textbf{Dev.} & \textbf{Route} \\
\midrule
$\alpha^{-1}$ & $3700/27=137.037$ & $+7.6$\,ppm &
  geometry from $T^4/\mathbb{Z}_3$ (Doc.~007), Galois confirmation (Doc.~338) \\
$\sin^2\theta_W$ (on-shell) & $2/9=0.22222$ & $-0.52\,\%$ &
  $A_5$ structure (Docs.~293, 372) \\
$\lambda_{\rm CKM}$ & $\xi^{1/6}=0.22602$ & $+0.46\,\%$ &
  Galois bridge (Doc.~336) \\
$v/E_P$ & $2.0389\times10^{-17}$ & $+1.10\,\%$ &
  $v$ from $\xi$ (R104), $E_P$ from our own chain (Doc.~375) \\
$E_{\rm bit}(L)$ & $\hbar c/L$ & --- &
  system-dependent, no universal bit value (Docs.~257, 326) \\
Dimension & $3{+}1$ from $T^4/\mathbb{Z}_3$ & --- &
  posit; $6=\mathbf{3}\oplus\mathbf{3'}$ as the bridge to HLV (Doc.~285) \\
Cosmic sector & --- & --- &
  bracketed (P39, Doc.~309) \\
\bottomrule
\end{tabular}
\end{center}

\section{What follows from this}

\begin{itemize}\itemsep2pt
\item Two frameworks are connected through more than one quantity: HLV through
  the $A_5$ structure and the dimension, OPH through the carrier, the scale
  and $v/E_P$. The rest are single comparisons.
\item The overlaps lie almost entirely in the lepton sector and in the
  geometry. The hadron sector appears in no comparison; that is not a
  selection but the state of the other sides.
\item $\alpha^{-1}$ and $v/E_P$ are currently the only quantities for which
  more than one framework states a number. They are therefore the only rows of
  a comparison table that data can tell apart.
\item The scale anchor is the same situation in a different place for every
  framework (Doc.~309); it is not a distinguishing feature.
\end{itemize}

\section*{Tool and method}
\addcontentsline{toc}{section}{Tool and method}
Claude (Anthropic) was used as a tool for: assembling the document list,
writing and running the check script, drafting the \LaTeX{} sources, and
maintaining the changelog. The tool computes and writes; the author decides
what is computed and written, judges every result and determines what it means
physically. The check script verifies that every cited document is present in
the corpus and recomputes every number of the second table from FFGFT
quantities.

\section{Result}

\begin{enumerate}
\item Eleven frameworks with comparison or bridge documents, two of them with
  several points of contact \KM.
\item Seven quantities where frameworks overlap, each with route and
  deviation \KM.
\item Only $\alpha^{-1}$ and $v/E_P$ are given a number by more than one
  framework.
\item The hadron sector appears in no comparison.
\end{enumerate}

\medskip
\noindent\textbf{Check script:} \texttt{2/python/Dok377\_Skripte/pruef\_377\_schnittstellen.py} (20/20)

\medskip
\noindent\textbf{Register (Doc.~190):} no entry.

\begin{thebibliography}{9}
\bibitem{hlv} Docs.~271, 272, 276, 282, 283, 285, 294, 297 (HLV).
\bibitem{oph} Docs.~364, 374, 375, 376 (OPH).
\bibitem{hyperbit} Docs.~324, 326, 336 (Hyperbit, $\mathbb{Z}_3\mathbb{C}$).
\bibitem{weitere} Docs.~049, 105, 251, 252, 269, 309, 345, 361.
\bibitem{quellen} Docs.~007, 013, 257, 293, 338, 372 (FFGFT sources of the numbers).
\end{thebibliography}
